\selectlanguage{english}

\textbf{\sffamily Conventions}

The totally-antisymmetric $ n $-dimensional Levi--Civita symbol $ \epsilon^{i_1 \dots i_n} $ is defined with the convention $ \epsilon^{1 \dots n} = +1 $.

Given $ \alpha \in \C^{n \times n} $, with $ n \in \N_0 $, its complex conjugate is denoted as $ \alpha^* $, its transpose as $ \alpha\tsp $ and its Hermitian conjugate as $ \alpha\dg \defeq (\alpha^*)\tsp $.

The Landau symbol for a function $ f : D \ssq \R \ra \C $ is defined by the condition $ \exists M \in \R \,:\, \abs{\smo(f(x))} \leq M \abs{f(x)} \,\,\forall x \in D $.

The imaginary unit is denoted by $ \img $, so that $ \C(\R) = \lspan_\R(1,\img) $.

\textbf{\sffamily Topology}

The empty set is denoted by $ \emptyset $ and the power set of a set $ A $ by $ \pwst{A} \defeq \{B : B \ssq A\} $. The counting numbers are $ \N \equiv \{1,2,3,\dots\} $, and the natural numbers are defined by $ \N_0 \equiv \{0\} \cup \N $.

The $ n $-dimensional sphere is denoted by $ \Sn{n} $, while the $ n $-dimensional disk by $ \Dn{n} $, so that $ \pa \Dn{n} = \Sn{n-1} $, where in general $ \pa \Sigma $ denotes the boundary of $ \Sigma $. The measure on $ \Sn{n} $ is denoted by $ \dd \Omega_n $, while its line element by $ \dd \Omega_n^2 $.

The permutation group of $ n $ objects, i.e. the $ n^\text{th} $ symmetric group, is denoted by $ \sy{n} $.

\textbf{\sffamily Multilinear Algebra}

Given two $ \K $-vector spaces $ V $ and $ W $, with $ \K $ a generic field, the space of all $ \K $-linear applications $ f : V \ra W $ is denoted by $ \Hom_\K(V,W) $: in particular, $ \Hom_\K(V , V) \equiv \End(V) $. The subset of $ \End(V) $ of all automorphisms of $ V $ is the automorphism group $ \Aut(V) $, which is a group under composition of morphisms.

Given a $ \K $-vector space $ V $, the space of all rank-$ (k,l) $ tensors on $ V $ is denoted by $ T^k_l(V) = \bigotimes^k (V) \otimes \bigotimes^l (V^*) $, the space of exterior $ p $-forms by $ \lv{p}(V) \ssq T^0_k(V) $ and the Grassmann algebra of $ V $ by $ \ga(V) = \bigoplus_{k = 0}^n \lv{k}(V) $.

\textbf{\sffamily Differential Geometry}

Given a fiber bundle $ \pi : E \ra M $, the space of smooth sections of $ E $ on $ M $ is denoted by $ \Gamma(E) $.

Given a smooth manifold $ \mathcal{M} $, the space of all smooth scalar functions on $ \mathcal{M} $ is denoted by $ \cm $, its tangent bundle by $ T\man $, its tangent rank-$ (k,l) $ tensor bundle by $ T^k_l \man = \bigsqcup_{p \in \man} T^k_l(T_p\man) $ and its alternating covariant $ k $-tensor bundle by $ \lm{k} = \bigsqcup_{p \in \man} \lv{k}(T_p\man) $. Then, the space of all smooth vector fields on $ \mathcal{M} $ is $ \xm = \Gamma(T\man) $, the space of all smooth $ (k,l) $-tensors on $ \man $ by $ \tns{k}{l}(\man) = \Gamma(T^k_l\man) $, the space of all smooth $ p $-forms on $ \mathcal{M} $ by $ \df{p} = \Gamma(\lm{k}) $ and the space of all smooth forms on $ \mathcal{M} $ by $ \dft \defeq \bigoplus_{k = 0}^n \df{k} $.

The exterior derivative is denoted by $ \dd $, the covariant exterior derivative by $ \dcov $, the partial derivative by $ \pa $, the nabla operator by $ \grad $ and the Laplace--Beltrami operator by $ \lap $. The codifferential is denote by $ \cdd $, and the Laplace--de Rham operator by $ \lapd = \dd \cdd + \cdd \dd $.

A list of ``important" Lie groups:

$ \GL{n,\K} \defeq \{\mt{A} \in \K^{n \times n} : \det \mt{A} \neq 0\} $ general linear group (Lie group for $ \K = \R , \C $)

$ \SL{n,\K} \defeq \{\mt{A} \in \GL{n,\K} : \det \mt{A} = 1\} $ special linear group (Lie group for $ \K = \R , \C $)

$ \On{n} \defeq \{\mt{A} \in \R^{n \times n} : \mt{A} \mt{A}\tsp = \mt{A}\tsp \mt{A} = \mt{I}_n\} $ orthogonal group

$ \SOn{n} \defeq \{\mt{A} \in \On{n} : \det \mt{A} = 1 $ special orthogonal group

$ \Un{n} \defeq \{\mt{A} \in \C^{n \times n} : \mt{A} \mt{A}\dg = \mt{A}\dg \mt{A} = \mt{I}_n\} $ unitary group

$ \SUn{n} \defeq \{\mt{A} \in \Un{n} : \det \mt{A} = 1 $ special unitary group

Given a Lie group $ G $, its associated Lie algebra is usually denoted by $ \mathfrak{g} $.










